% \m@th in a formula, which MathJax does not know. hologo's BibTeX logo
% sets "ib" in a smaller font that starts with an empty $\m@th$: in text,
% in a heading and its entry in the table of contents, in italic, in a
% footnote, in a theorem, in sans-serif, typewriter, and bold text, as
% \Hologo, through BibTeX8, and as the sf variant. A copy of the same
% construct in a document's macro, \m@th before content in \mbox and
% \ensuremath, and in the body of a math macro. Guards: the sc variant,
% which has no math, and logos that run \m@th outside their formula.
% test-hologo-plain has the same logo without hyperref.
\documentclass{article}
\usepackage{amsmath,amsthm}
\usepackage{hologo}
\usepackage{hyperref}
\newtheorem{theorem}{Theorem}
\makeatletter
\newcommand\OldLaTeX{L\kern-.36em{\setbox0\hbox{T}\vbox to\ht0{\hbox{$\m@th$\csname S@\f@size\endcsname\fontsize\sf@size\z@\math@fontsfalse\selectfont A}\vss}}\kern-.15em\TeX}
\newcommand\Boxed{\mbox{$\m@th q+r$}}
\newcommand\Ensured{\ensuremath{\m@th\alpha\beta}}
\newcommand{\nrm}[1]{\m@th\lVert #1\rVert}
\makeatother
\begin{document}
\tableofcontents
\section{Using \hologo{BibTeX} files}
Text: \hologo{BibTeX} now.

Italic: \emph{\hologo{BibTeX} here}.

Footnote: see below.\footnote{Notes: \hologo{BibTeX} there.}

\begin{theorem}
Theorem: \hologo{BibTeX} holds.
\end{theorem}

Sans: \textsf{\hologo{BibTeX}}. Typewriter: \texttt{\hologo{BibTeX}}. Bold: \textbf{\hologo{BibTeX}}.

Capital: \Hologo{BibTeX}. Eight: \hologo{BibTeX8}. Variant sf: \hologoVariant{BibTeX}{sf}.

Copy: \OldLaTeX.

Boxed: \Boxed. Ensured: \Ensured. Norm: $\nrm{w}$.

Guards: \hologoVariant{BibTeX}{sc}, \hologo{LaTeXe}, \hologo{eTeX}, \hologo{ExTeX}, and \hologo{TTH}.
{\hologoLogoSetup{BibTeX}{variant=sc}Setup: \hologo{BibTeX}.}
\end{document}
